\documentclass[10pt,a4paper]{amsart}
\usepackage[margin=19mm]{geometry}
\usepackage{amsmath,amssymb,mathtools}
\usepackage[scr=rsfs,cal=euler]{mathalfa}
\usepackage{xfrac}
\usepackage[unicode,hidelinks,bookmarks=false]{hyperref}
\newcommand{\ie}{i.e.\ }
\newcommand{\cf}{cf.\ }
\newcommand{\resp}{respectively\ }
\newcommand{\Subsection}{\subsection}
\providecommand{\nobreakdash}{\nobreak\hbox{-}}
\setlength{\emergencystretch}{2em}
\multlinegap0pt
\usepackage{amssymb}
\usepackage{tikz}
\usepackage{enumerate}
\usepackage[shortlabels]{enumitem}
\usepackage{mathtools}
\newtheorem{theorem}{Theorem}
\newcommand{\R}{\mathbb{R}}
\title{Upper Bounds on Covering Minima of Convex Bodies}
\author{Katarina Krivoku\'ca}
\date{Source: arXiv:2601.15173v1 (CC BY 4.0)}
\begin{document}
\maketitle

\noindent\emph{Source and licence.} Upper Bounds on Covering Minima of Convex Bodies. Version arXiv:2601.15173v1 (CC BY 4.0), distributed by arXiv under the Creative Commons Attribution 4.0 International licence, http://creativecommons.org/licenses/by/4.0/, as recorded in the arXiv licence metadata for this exact version and shown on the abstract page https://arxiv.org/abs/2601.15173v1. Full licence notice: \texttt{licenses/arxiv-2601.15173-CC-BY-4.0.txt}.\par\smallskip

\noindent\textit{Editorial scope.} Counted unit: the article's theorem proj (source0.tex lines 177--182). The article's complete original proof of it is delivered with it (source0.tex lines 240--267, one \texttt{proof} environment of 28 lines, which the article places away from the statement). No mathematics is written, repaired or re-derived: the statement and the proof are the article's own lines, in the article's own notation, and no retained line is substituted or re-flowed.  No further source-local context is needed: every cross-reference of every retained line resolves inside the two delivered spans.  Nothing was omitted from any retained span, so the package carries no omission record.  No retained line contains a citation, so the wrapper carries no bibliography and no bibliography entry.  No retained line contains a figure environment, an image-inclusion command or drawing code, so no image is delivered and the package declares no asset dependency.  Editorial formatting only, in addition: an AMS \texttt{amsart} preamble that restates the article's own declarations and macros used by the retained lines, namely the environments \texttt{theorem} on separate counters, together with the standard packages those lines need.  The retained environments print in this document as Theorem 1 in order of appearance, whereas the article prints the same items with the numbering of its own full document.  The complete native source of the article as distributed in the arXiv e-print for this exact version is bundled unchanged as \texttt{sources/193137/source0.tex}.
\begin{theorem}\label{thm: ub proj}
Let $K\subseteq\R^d$ be a convex body containing the origin, $\Lambda\subseteq \R^d$ a lattice and $V\subseteq \R^d$ a rational linear subspace, with $\dim V=\ell$ and $i\in[d]$. If $\pi_V$ denotes the orthogonal projection of $\R^d$ to $V$, the following inequality holds:
\begin{align}\label{eq:ub proj}
\mu_i(K, \Lambda)\leqslant \max_{\substack{0 \le j \le \ell \\ 0 \le i-j \le d-\ell}} \mu_j(\pi_V(K), \pi_V(\Lambda))+ \mu_{i-j}(K\cap V^\perp, \Lambda \cap V^\perp)
\end{align}
\end{theorem}

\begin{proof}[Proof of Theorem \ref{thm: ub proj}]
Let $x+U$ be an arbitrary $(d-i)$-dimensional affine subspace of $\R^d$, where $x\in \R^d$ and $U\subseteq \R^d$ is a linear subspace. It suffices to show that $x+U$ intersects $(\mu_j(\pi_V(K), \pi_V(\Lambda))+ \mu_{i-j}(K\cap V^\perp, \Lambda \cap V^\perp))K + \Lambda $, for some $0 \le j \le \ell,\; 0 \le i-j \le d-\ell$.

Let $\dim(\pi_V(U))=\ell-j$ and $\dim(U\cap V^\perp)=(d-i)-(\ell-j)=(d-\ell)-(i-j)$. For brievity, let $\mu_j:=\mu_j(\pi_V(K),\pi_V(\Lambda))$ and $\mu_{i-j}:=\mu_{i-j}(K\cap V^\perp, \Lambda \cap V^\perp)$.

Since $\pi_V(x)+\pi_V(U)$ is a $(\ell-j)$-dimensional affine subspace of $V$, there exist $u_1\in U$, $p\in K$ and $a\in \Lambda$ such that 
\begin{align}\label{eq:cmj}
\pi_V(x)+\pi_V(u_1)=\mu_j \pi_V(p)+\pi_V(a)
\end{align}

Denote by $y=\pi_{V^\perp}(x)+\pi_{V^\perp}(u_1)-\mu_j\cdot \pi_{V^\perp}(p)-\pi_{V^\perp}(a)\in V^\perp$.


\noindent Notice that $y\in V^\perp$ implies that $\dim((y+U)\cap V^\perp)=\dim(U\cap V^\perp)=(d-\ell)-(i-j)$. 

\noindent Since $(y+U)\cap V^\perp$ is a $((d-\ell)-(i-j))$-dimensional affine subspace of $V^\perp$, there exist $u_2 \in U\cap V^\perp$, $q\in K\cap V^\perp$ and $b\in \Lambda \cap V^\perp$ such that

$$y+u_2=\mu_{i-j}q+b$$
$$\Rightarrow \pi_{V^\perp}(x)+\pi_{V^\perp}(u_1)-\mu_j\cdot \pi_{V^\perp}(p)-\pi_{V^\perp}(a)+u_2=\mu_{i-j}\cdot q+b$$
\begin{align}\label{eq:cmi-j}
\Rightarrow \pi_{V^\perp}(x)+\pi_{V^\perp}(u_1)=\mu_j\cdot \pi_{V^\perp}(p)+\mu_{i-j}\cdot q+\pi_{V^\perp}(a)+b
\end{align}
Now, adding the equations \ref{eq:cmj} and \ref{eq:cmi-j} we get:

$$ x+u_1+u_2=\mu_j\cdot p+\mu_{i-j}\cdot q+a+b$$
$$\Rightarrow x+u_1+u_2=(\mu_j+\mu_{i-j})\left(\frac{\mu_j}{\mu_j+\mu_{i-j}}p+\frac{\mu_{i-j}}{\mu{j}+\mu_{i-j}}q\right)+a+b.$$

This point verifies that $x+U$ intersects $(\mu_j+\mu_{i-j})K + \Lambda$. Since $j$ depends on $U$, and by construction $j$ can take any value for which $0 \le j \le \ell,\; 0 \le i-j \le d-\ell$, it follows that $\mu_i(K,\Lambda)$ is bounded from above by the maximum of $\mu_j+\mu_{i-j}$ for all $j$ in this index set. \end{proof}

\end{document}
